% Sommario (abstract) dell'allegato: solo il testo, da racchiudere in \begin{abstract}...\end{abstract} o in un paragrafo «Sommario».
Si descrive un prototipo di percezione che rileva elementi mobili e ostacoli con una camera di profondità a stereo attivo, frontale e ferma, e che continua a dare un'indicazione dichiarata quando un ostacolo nasconde il soggetto. Il sistema tiene distinte tre informazioni: l'osservato; il non osservato ma previsto, con una regione e un livello; il non noto. Lo stato «libero» è scritto solo per il suolo osservato nel fotogramma corrente; l'avviso di possibile presenza non termina senza osservazione diretta. Nelle ipotesi dichiarate sono dimostrati che ciò che non è osservato non è mai «libero», il contenimento del soggetto nascosto nell'insieme di possibile presenza (per una lettura esatta del suolo, senza il margine che il sistema usa) e la condizione sulla fascia di guardia. In simulazione, con modello ideale del sensore e un soggetto alla volta, la regione di previsione è emessa a ogni ciclo dell'occlusione una volta riconosciuto l'ingresso nell'ombra (0 interruzioni dopo l'ingresso su $23\,402$ passaggi nascosti; in 928 di essi la traccia si chiude prima dell'ombra e non c'è regione, e con una regola aggiuntiva, attiva solo dopo una prova d'ingresso, 24 restano senza avviso); con un ostacolo fermo largo $2\unita{m}$, il bordo dell'ombra noto e il soggetto sul proprio percorso, la regione contiene il soggetto in almeno 9 casi su 10 fino a $2\unita{s}$ e in almeno 85 su 100 fino a $3\unita{s}$; con un ostacolo largo $1\unita{m}$ la soglia di 9 su 10 non è raggiunta, e su profondità simulata dalle schede tecniche questa regola non è stata eseguita. Dietro un ostacolo fermo l'avviso resta acceso, e nella realizzazione provata resta acceso anche dopo la rimozione dell'ostacolo. Il classificatore degli oggetti è un componente richiesto dal sistema: in questo lavoro non è realizzato né valutato, ed è simulato con etichette assegnate a mano. Nessuna prova è stata eseguita su un sensore fisico: tutti i risultati sono ottenuti a camera ferma, al chiuso, fino a $3\unita{m}$, e non sono stati verificati fuori da queste condizioni.
